\begin{center}
{\Large\bfseries The $(3,4,\infty)$ family of complex structures on $S^6$:\\periods, moduli and the ends of the disc\par}
\medskip
{\small AI-produced research\quad 17 November 2026\par}
\end{center}
\begin{abstract}
For the standard marked $(3,4,\infty)$ family, the admissible parameter is
$0<|w|<216$, where $w=\exp(-\pi i c)$, and $X_c\simeq X_{c'}$ exactly when
$c'-c\in2\mathbb Z$. We determine its normalized periods, sharp norm bound,
integral monodromy, cohomology, real forms, special fibres and arithmetic
specializations. The specified affine family has no proper analytic filling
at zero after any finite parameter cover; the circle $|w|=216$ is a natural
boundary for proper flat extensions. The construction is credited to Claude
and Alp\"oge and to Engel; the standard classification, tangent cohomology,
Kuranishi germ and automorphisms are credited to Du, Guo, Kusakabe and Xie,
version 1. Complete independent derivations accompany the period, integral,
analytic and arithmetic calculations. This writing and its proof comparisons
are AI-produced by GPT-6 Astra (OpenAI, ChatGPT Pro mode, effort Pro),
17 November 2026; the verification record distinguishes complete readings,
source checks and rerun certificates. No human has checked the mathematics.
A literal global coefficient atlas and an unconditional transcendence claim
remain in a separate final part.
\end{abstract}

\section{Result}
\subsection{Statement}
We use the marked construction and orientation conventions of
Section~\ref{def:common-conventions}. The standard family has the parameter
domain and equivalence relation
\[
 \Im c<\frac{3}{\pi}\log6,\qquad
 w=e^{-\pi i c},\qquad 0<|w|<216,\qquad
 X_c\simeq X_{c'}\ \Longleftrightarrow\ c'-c\in2\mathbb Z.
\]
The norm maximum is $\sup D_{\mathrm{part}}=-(3/\pi)\log6$, attained only
at the order-four base point. The monodromy group is the explicit index-twelve
kernel of Theorem~\ref{fam:src:F.1.1}. The complete Hodge and Bott--Chern
tables are in \cref{inv:A3.1,inv:C2.1}; in particular
$h^{1,1}=h^{2,2}=2$ and $h^{1,2}=h^{2,1}=1$.
The tangent dimensions are $(1,1,1,0)$, the reduced Kuranishi germ is a
smooth line, and $\Aut(X_c)=\mathbb C^*$.
The no-filling and boundary assertions, with their exact hypotheses, are
\cref{push:src:B.4.33,push:src:D.1.3}.
For instance $c=0$ gives $w=1$ in the parameter disc, whereas $c+2$ gives
the same biholomorphism class.

That this family contradicts \citet[Corollary~2.3]{CDP20} is due to
\citet[Section~1.2 and Theorem~10.5]{AC2026}. Paper~II, Propositions~5.1 and 13.2, gives the general analysis and replacement criteria; Paper~III, Theorems~2.4 and 6.52 and Corollary~7.32, gives the map, framing and manifold results. No conditional arithmetic statement is used
to establish the complex family or its parameter disc.

\subsection{Status}
\begin{tabular}{@{}p{30mm}p{121mm}@{}}
\toprule
Kind & Theorems, computations, corrections, counterexamples and credited re-derivations.\\
Claim & The body contains checked proofs; unresolved questions form the final part.\\
Checks & Complete readings and re-derivations, source-theorem checks, exact and interval reruns; the per-result concordance gives their scopes.\\
Human understanding & None claimed; no human has checked the mathematics.\\
Formalization & None; no proof-assistant certification is claimed.\\
Open points & Two questions, with their complete partial arguments and precise missing steps.\\
\bottomrule
\end{tabular}

\subsection{Known and new}
\begin{tabular}{@{}p{41mm}p{70mm}p{38mm}@{}}
\toprule
Statement & Source and exact scope & Present treatment\\
\midrule
Construction, topology, canonical pencil &
\citet[Main Theorem; Sections~3--7; Theorem~9.1]{AC2026};
\citet[main theorem]{Engel2026} & Full construction and independent calculations.\\
Standard classification &
\citet[Theorem~1.4]{DuEtAl2026v1} & Full period-shift proof in the appendices.\\
Tangent and Kuranishi statements; automorphisms &
\citet[Theorem~1.4; Theorem~6.12; Proposition~6.16]{DuEtAl2026v1} & Re-derived after reading that version.\\
Structure-sheaf cohomology; Picard restriction &
\citet[Lemma~9.12; Proposition~9.13; Theorem~10.5]{AC2026} & Direct-image, residue and period derivations retained.\\
Contradiction to CDP20 &
\citet[Section~1.2; Theorem~10.5]{AC2026} & Credited here; general analysis in Paper~II.\\
Period bounds, discrete refinements, full invariant calculations and ends &
The exact comparison and correction entries in Appendix~\ref{app:I:result-index} & Statements matched individually; distinct proofs retained.\\
\bottomrule
\end{tabular}

The standard classification is Theorem~1.4 of version~1
(23 September 2026); version~2 does not contain it. Our own proof was in the
public working edition by 8 September 2026. We print that complete proof
in the present notation. The tangent cohomology, Kuranishi germ and full
automorphism group are explicitly re-derived results.
The literature comparison uses the original Claude--Alp\"oge PDF, Engel,
both Du--Guo--Kusakabe--Xie versions and the additional sources at their
recorded locators. The original PDF controls discrepancies with its
searchable transcription.

\enlargethispage{\baselineskip}
\subsection{Proof map}
The root charts, affine origins and cusp fan fix the actual compact family.
The period lattice gives the finite norm obstructions and the even-shift
classification. Theta identities, analytic inequalities and separately
rerun interval bounds determine the disc. Root-curve direct images and the
logarithmic cotangent sequence give the cohomology and deformation results.
At zero the obstruction is the radial class of the actual affine origins;
at the outer circle the intrinsic multiple fibre has an Albanese invariant
with a natural boundary. Arithmetic proofs retain all hypotheses and error
estimates. The construction and credited re-derivations appear in the
appendices; Appendix~\ref{app:I:result-index} gives the complete source-label,
proof and verification concordance.

\clearpage
\section{Common conventions and the marked construction}
\label{def:common-conventions}
